\documentclass[11pt]{article}
\usepackage[margin=1in]{geometry}
\usepackage{hyperref}
\title{LSTM Forecasting for Mauna Loa CO$_2$}
\author{Devis W. Saputra}
\begin{document}
\maketitle

\begin{abstract}
I use 24 weeks of Mauna Loa atmospheric CO$_2$ history to forecast the next weekly value with a small LSTM. The evaluation is chronological and normalization is fitted on the training period only.
\end{abstract}

\section{Question}
Can a 24-week historical window support useful one-step-ahead forecasting with a compact LSTM?

\section{Data and Method}
The CO$_2$ series is resampled to weekly means and missing weeks are interpolated. I create 2,260 windows, use the first 1,808 for training, and hold out the final 452. The LSTM has 32 hidden units and a one-value linear output.

\section{Training}
The model trains for 20 epochs with Adam at learning rate 0.005 and mean-squared error loss.

\section{Results}
After fitting normalization on the training period only, the recorded hold-out RMSE is 4.4853 and MAE is 4.3378.

\section{Limitations}
The evaluation uses one chronological split and does not yet include naive or autoregressive baselines. Those comparisons are needed before judging the LSTM's practical value.

\section{Reproduction}
Install \texttt{requirements.txt} and run \texttt{python src/run\_experiment.py}.

\bibliographystyle{plain}
\bibliography{references}
\end{document}
